\subsection{\texorpdfstring{Complete subsets of \(\mathcal{L}_{\mathfrak{S}}(\mathbf{E};\mathbf{F})\)}{Complete subsets of \textbackslash mathcal\{L\}\_\{\textbackslash mathfrak\{S\}\}(\textbackslash mathbf\{E\};\textbackslash mathbf\{F\})}}

PROPOSITION 11. --- Let E and F be two locally convex spaces, S a cover of E consisting of bounded subsets. If F is Hausdorff and quasi-complete (III, p.~8), then every equicontinuous subset H of \(\mathcal{L}(\mathrm{E};\mathrm{F})\) which is closed for the S-topology is a complete uniform subspace of \(\mathcal{L}_{\mathfrak{S}}(\mathrm{E};\mathrm{F})\) .

Since H is bounded in \(\mathcal{L}_{\mathfrak{S}}(\mathrm{E};\mathrm{F})\) (III, p.~22, prop. 9) and closed in \(F^{E}\) for the S-topology (III, p.~16, prop. 4), this follows from cor. 3 of GT, X, § 1, No.~5.

Remark 1. --- Let M be a complete uniform subspace of \(\mathcal{L}_{\mathfrak{S}}(\mathrm{E};\mathrm{F})\) . For every set of bounded subsets \(\mathfrak{S}' \supset \mathfrak{S}\) of E, the \(\mathfrak{S}'\) -topology is finer than the \(\mathfrak{S}\) -topology on \(\mathcal{L}(\mathrm{E};\mathrm{F})\) ; on the other hand, there exists a fundamental system of neighbourhoods of 0 for the \(\mathfrak{S}'\) -topology which are closed for the topology of simple convergence (III, p.~13, Remark 2), and a fortiori for the \(\mathfrak{S}\) -topology. We conclude (GT, III, § 3, No.~5, cor. 1) that M is complete for the \(\mathfrak{S}'\) -topology.

COROLLARY. --- Let E and F be two locally convex spaces, H an equicontinuous subset of \(\mathcal{L}(\mathrm{E};\mathrm{F})\) . If F is Hausdorff and quasi-complete and if a filter \(\Phi\) on H converges simply at all points of a total subset T of E, then there exists a continuous linear mapping u from E into F such that \(\Phi\) converges uniformly to u on every precompact subset of E.

For, by virtue of prop. 5 (III, p.~17) \(\Phi\) is a Cauchy filter for the uniform structure of precompact convergence in E; by prop. 11, the closure \(\overline{\mathrm{H}}\) of H in \(\mathcal{L}_{pc}(\mathrm{E};\mathrm{F})\) is complete and so \(\Phi\) converges uniformly on every precompact subset of E to a mapping \(u\in \overline{\mathrm{H}}\) .

Remark 2. --- Let \((u_{n})\) be a sequence of continuous linear mappings from a Banach space E into a Banach space F; it may happen that \((u_{n}(x))\) has a limit at every point of an everywhere dense vector subspace T of E, without the sequence \((u_{n})\) being bounded in the normed space \(\mathcal{L}(\mathrm{E};\mathrm{F})\) . For example, take E to be the space of all continuous numerical functions on R, tending to zero at infinity, with the norm \(\|f\| = \sup_{x \in \mathbf{R}} |f(x)|\) and let T be the subspace of continuous numerical functions

with compact support. The sequence of continuous linear mappings \(f \mapsto nf(n)\) from E into R converges to 0 for all \(f \in T\) , but is not bounded in \(\mathcal{L}_{b}(E; \mathbf{R})\) . The same example shows that in the space \(\mathcal{L}(T; \mathbf{R})\) , a sequence \((v_{n})\) may be simply convergent and non-bounded for the topology of bounded convergence.

On the other hand, the sequence of continuous linear mappings \(f \mapsto \sum_{k=1}^{n} f(k)\)

is a Cauchy sequence in \(\mathcal{L}(\mathrm{T};\mathbf{R})\) for the topology of simple convergence, but does not tend to a limit in \(\mathcal{L}(\mathrm{T};\mathbf{R})\) for this topology.

PROPOSITION 12. --- Let E be a bornological locally convex space, F a complete locally convex Hausdorff space and S a family of bounded subsets of E containing the image of every sequence converging to 0. Then the space \(\mathcal{L}_{\mathfrak{S}}(\mathrm{E};\mathrm{F})\) is complete.

Let \(\Phi\) be a Cauchy filter in \(\mathcal{L}_{\mathfrak{S}}(\mathrm{E};\mathrm{F})\) . Then \(\Phi\) is a Cauchy filter for the topology of simple convergence, hence converges in \(\mathrm{F}^{\mathrm{E}}\) ; moreover, its limit \(u\) is a linear mapping from E into F and \(\Phi\) converges to \(u\) uniformly on every set of \(\mathfrak{S}\) (GT, X, § 1, No.~5, prop. 5). It follows that the image under \(u\) of a sequence converging to zero is a sequence converging to zero, hence, that \(u\) is continuous, since E is bornological (III, p.~11, prop. 1, (iii)).

COROLLARY 1. --- The strong dual of a bornological space is complete.

COROLLARY 2. --- Let E be a semi-normed space, and F a Banach (resp. Fréchet) space. The space \(\mathcal{L}_{b}(\mathrm{E};\mathrm{F})\) is a Banach (resp. Fréchet) space. In particular, the dual of a semi-normed space is a Banach space.

